% Indexing and Bibliography Citations Example
\documentclass[12pt]{article}
\usepackage{makeidx}
\makeindex

\begin{document}

\section*{Error Troubleshooting Strategy}
When LaTeX encounters an error, it provides a line number and a descriptive error message in the log file (e.g., \texttt{Undefined control sequence} or \texttt{Too many columns in tabular}). 
\index{Debugging!Error Logs}

\section*{Subject Indexing}
Important terms in large documents can be indexed automatically using the makeidx package. For example, we can index LaTeX commands\index{LaTeX!Commands} and compilation steps\index{Compilation}.

\section*{References and Citations}
Academic and technical documents rely on BibTeX for reference management. A typical citation looks like this: as explained by Lamport \cite{lamport}.
\index{Bibliographies!BibTeX}

\begin{thebibliography}{9}
\bibitem{lamport}
Leslie Lamport,
\emph{LaTeX: A Document Preparation System},
Addison-Wesley, 1994.
\end{thebibliography}

\printindex

\end{document}